% ==============================================================
%  Chapter 3: AI for DevOps Engineers
% ==============================================================
\chapter{AI for DevOps Engineers}
\label{ch:devops}

DevOps engineers deal with an unusually wide surface area: infrastructure code,
pipelines, container orchestration, security policies, monitoring configuration, and
incident response --- often at the same time.
AI is a genuine force multiplier here because most of the work involves translating
intent into configuration syntax, a task AI handles extremely well.

\section{Infrastructure as Code}
\label{sec:iac}

Writing Terraform, Ansible, Helm, and Pulumi configurations is one of the clearest
wins for AI assistance.
The configuration formats are well-documented, schema-constrained, and follow patterns
the model has seen in millions of public repositories.

\textbf{What works well:}
\begin{itemize}
  \item Generating a complete Terraform module from a requirements description
  \item Converting between IaC formats (e.g., CloudFormation to Terraform)
  \item Adding lifecycle rules, tags, and provider constraints you forgot
  \item Explaining what an existing Terraform resource will actually provision
  \item Writing Ansible playbooks for standard system configuration tasks
\end{itemize}

\textbf{Effective prompts:}
\begin{itemize}
  \item \textit{``Write a Terraform module for an AWS ECS service. Requirements: Fargate, 512 CPU, 1 GB memory, private subnets, ECR image, environment variables from SSM. Include an IAM execution role.''}
  \item \textit{``Review this Terraform configuration for security issues: over-permissive IAM, public S3 buckets, unencrypted EBS.''}
  \item \textit{``Convert this Docker Compose file to a Kubernetes Deployment and Service manifest.''}
\end{itemize}

\begin{warnbox}
AI does not know your cloud account state, existing resources, or organisation-specific
naming conventions.
Always run \texttt{terraform plan} before \texttt{apply} and review the diff carefully.
AI-generated IaC often uses permissive defaults that need tightening for production.
\end{warnbox}

\section{CI/CD Pipelines}
\label{sec:cicd}

Pipeline YAML for GitHub Actions, GitLab CI, and Jenkins follows strict schemas that
AI handles well.
Common use cases:
\begin{itemize}
  \item Generating a complete pipeline from a description of your workflow (build, test, scan, deploy, rollback)
  \item Adding matrix builds, caching strategies, and environment-specific deployments
  \item Translating a Jenkins pipeline to GitHub Actions
  \item Diagnosing why a pipeline step is failing from logs
\end{itemize}

\begin{tipbox}[Pipeline Prompt Pattern]
\textit{``Write a GitHub Actions workflow that: (1) triggers on push to main and PRs, (2) runs pytest with coverage on Python 3.12 and 3.13, (3) builds and pushes a Docker image to GHCR on main only, (4) deploys to staging via SSH after a successful push. Use caching for pip dependencies.''}
\end{tipbox}

\textbf{Diagnosing failed pipelines:}
Paste the failing step's full log output and ask:
\textit{``What is causing this CI failure and what is the fix?''}
AI is excellent at reading pipeline logs and identifying root causes --- permission
errors, missing environment variables, incorrect paths, version mismatches.

\section{Containers and Kubernetes}
\label{sec:k8s}

\textbf{Dockerfile optimisation:}
Paste your current Dockerfile and ask:
\begin{itemize}
  \item \textit{``Optimise this Dockerfile for image size. Explain each change.''}
  \item \textit{``Reorder these layers to maximise cache hit rate.''}
  \item \textit{``What security issues exist in this Dockerfile (running as root, secrets in ENV, etc.)?''}
\end{itemize}

\textbf{Kubernetes manifests:}
AI generates Deployment, Service, ConfigMap, HPA, and PodDisruptionBudget manifests
reliably.
Useful for generating base manifests, then customising for your environment.

\begin{tipbox}
Ask AI to include resource requests and limits, liveness/readiness probes, and
non-root user configuration by default.
These are the settings most often forgotten and most often flagged in security audits.
\end{tipbox}

\section{Monitoring, Alerting, and Incident Response}
\label{sec:monitoring}

\textbf{What AI does well:}
\begin{itemize}
  \item Generating Prometheus alert rules from verbal descriptions
    (\textit{``Write an alert that fires if p99 latency exceeds 500 ms for 5 minutes''})
  \item Writing PromQL and LogQL queries when you describe what you want to measure
  \item Drafting runbooks and incident response playbooks
  \item Explaining what a specific alert or metric means in plain language
\end{itemize}

\textbf{During incidents:}
Paste your log output and metrics anomaly and ask:
\textit{``What does this pattern suggest about root cause? What would you check next?''}
AI cannot see your live systems, but it can reason about log patterns and suggest a
diagnostic sequence quickly when you are under pressure.

\begin{warnbox}
AI cannot access your live monitoring dashboards, alert history, or deployment state.
It works from what you paste. Be precise and complete when describing an incident.
\end{warnbox}

\section{Security and Compliance Automation}
\label{sec:sec-compliance}

\textbf{Useful security tasks for AI:}
\begin{itemize}
  \item Reviewing IAM policies for over-permission (\textit{``Does this IAM policy follow least-privilege? What can be removed?''})
  \item Generating security group rules and network policies from verbal requirements
  \item Writing Semgrep or OPA/Rego rules for custom policy enforcement
  \item Drafting threat model sections for a new service
  \item Explaining CVEs and assessing their relevance to your stack
\end{itemize}

\begin{tipbox}
For cloud security reviews, paste your policy JSON and ask AI to identify issues against
the CIS benchmark or AWS Well-Architected Framework security pillar.
This is faster than reading the benchmark yourself and catches 80\% of common issues.
\end{tipbox}
